% Source: source_md/intelligence_part2_combined_ja.md
\section{アニミズム — $L_E$ を環境直接相互干渉で維持する}


\S{}5 で創発停止層 $L_E$ と、その個体間ネットワーク $\Phi_{\text{net}}$ を確立した。本章から \S{}9 まで、$L_E$ と $\Phi_{\text{net}}$ が歴史的・社会的に取る形態を分析する。本章はその最初の、最も基底的な形態であるアニミズムを扱う。

あらかじめ \S{}0.4 の非目標を再確認する。本章はアニミズムの擁護でも批判でもなく(NG1)、特定の歴史的・文化的伝統の解釈でもなく(NG5)、対象に認知や意図を帰属するものでもない(NG6)。扱うのは、アニミズムという停止構造の動力学的形態のみである。

\bigskip


\subsection{アニミズムの構造的定義}


アニミズムを、信念体系の内容によってではなく、停止構造の形態によって定義する。

\textbf{定義 6.1 (構造としてのアニミズム)}

\begin{quote}
アニミズムは、沈黙的世界に応答チャネルを注入する操作である。$\nabla C$ が一方向のみで応答ループが閉じない領域(沈黙領域)に対して、応答性を局所的に付与し、その領域を応答的に扱えるようにする動力学的操作を指す。
\end{quote}


\S{}4 で示した通り、沈黙的世界では推論が無限後退に陥る。環境が応答を返さないため、推論を停止させる外部応答が来ない。アニミズムは、この沈黙領域に応答チャネルを注入することで、無限後退を回避する。

C-幾何学の語彙で述べれば、アニミズムは応答チャネル密度 $\rho_{\text{resp}}$ を沈黙領域において局所的に増大させる操作である。山・川・風・石といった、本来は推論主体に応答を返さない対象に対して、応答性を付与する。これにより、$H_{ji} \approx 0$ だった領域が、$H_{ji} > 0$ として扱われるようになる。

\bigskip


\subsection{事実性ではなく構造的変換}


ここで決定的な区別を置く。アニミズムが行うのは、世界についての事実の主張ではなく、停止構造の変換である。

「山が怒る」「川が許す」「風が意図する」といった表現は、山や川や風が実際に感情や意図を持つという事実主張として読まれるべきではない。本論文の枠組みでは、これらは沈黙領域に応答チャネルを注入する構造的操作の表現である。重要なのは、山が本当に怒るかどうかではなく、山を応答的に扱うことによって推論主体の停止構造がどう変わるか、である。

応答的に扱われた対象は、推論主体に対して応答を返す位置に置かれる。主体は対象に問いかけ、対象からの応答(と解釈されるもの)を受け取り、それによって推論を停止できる。停止が可能になることが要点であって、対象の内的状態の有無は問題ではない。

これは \S{}5.6 で論じた投影の一形態である。応答性という影響力が、沈黙対象に投影される。意識編の projection の定義——ある対象が不均等な影響力を獲得する操作——がここでも作用している。ただし \S{}5.6 と同じく、この投影は主体の内的状態の対象への帰属(心理学的 projection)とは区別され、停止構造の動力学的変換として記述される。

\bigskip


\subsection{因果推論の否定ではなく保留}


アニミズムはしばしば、因果的・科学的説明の対立物と見なされる。本論文の枠組みでは、この見方は正確ではない。

アニミズムは因果推論を否定しない。それは因果推論を保留する。「なぜ災害が起きたのか」という問いに対して、アニミズムは「自然の意図」という応答チャネルを注入することで、いったん停止を可能にする。しかしこれは、因果的探究を永久に禁じるものではない。それは説明の遅延であって、説明の拒否ではない。

この区別は重要である。説明拒否であれば、因果推論は不可能になる。しかし説明遅延であれば、因果推論は可能なまま保たれ、ただ停止点が外部化されることで無限後退が回避される。アニミズムは推論を止めるのではなく、推論を操作可能な範囲に保ちながら、停止を供給する。

これは \S{}8(宗教)で詳述する「因果推論の部分的解放」の最も原初的な形態である。

\bigskip


\subsection{低コスト停止装置}


アニミズムが歴史的に広く現れた理由を、動力学的コストの観点から説明できる。

アニミズムは低コストの停止装置である。沈黙領域に応答チャネルを注入するのに、複雑な計算も、体系的な検証も、制度的基盤も要しない。最小限の構造的操作——対象を応答的に扱うこと——だけで機能する。

この低コスト性が、アニミズムの普遍性を説明する。停止を供給する他の形態(後述する神・宗教・社会)は、より高度な構造を要する。アニミズムはそれらに先立つ、最も少ない構造的前提で成立する停止形態である。これは優劣の問題ではなく、構造的コストの問題である。

\bigskip


\subsection{責任の所在と行動可能性}


応答的に扱われる世界は、推論主体に対して責任の場を提供する。

沈黙的世界では、出来事は誰のものでもなく、応答も帰責もない。これに対し、応答的に扱われた世界では、出来事は応答可能な対象に帰される。これにより、主体は出来事に対して行動可能になる——働きかけ、応答を受け取り、関係を調整する余地が生まれる。

この行動可能性は、停止構造の副産物である。応答チャネルが注入されることで停止が可能になると同時に、その応答チャネルを通じた働きかけが可能になる。停止と行動可能性は、同じ応答チャネル注入の二つの帰結である。

繰り返すが、これは応答対象が実際に責任主体であることを意味しない。責任の場が構造的に成立することが、主体の行動可能性を支える、という動力学的記述である。

\bigskip


\subsection{三層動力学における位置 — 環境直接相互干渉型 $L_E$}


アニミズムを三層動力学の語彙で位置づける。

アニミズムは、創発停止層 $L_E$ を環境との直接相互干渉によって維持する形態である。\S{}5.4 で見た粘菌のプロファイル——環境相互干渉によって $L_E$ を自律創発させる——と構造的に同位置にある。

具体的には、アニミズムにおける $\Phi_{\text{net}}$ は環境的・分散的である。停止を委ねる先(応答チャネルの注入先)が、個々の自然物・自然現象に分散している。山には山の、川には川の応答チャネルがあり、$L_E$ ノードが環境全体に分散して配置される。これは後述する神(\S{}7)の中心化された $\Phi_{\text{net}}$ とは対照的である。

そしてアニミズムは、閾値移動能力 $\mu_\theta$ の最初の社会的具現化である。\S{}2.7 で論じた通り、$\mu_\theta$ は個体内で運動可能領域を組み替え、停止演算子設置領域を投影する能力である。アニミズムは、この投影を文化的装置として共有・安定化させる。個体が単独で行う閾値移動を、共有された応答チャネルの体系として外部化する。これにより、$\mu_\theta$ の作用が個体を超えて persistent になる。

\subsubsection{粘菌との構造的同位について}


\S{}5.4 で予告した通り、粘菌の環境相互干渉型 $L_E$ とアニミズムの $L_E$ 維持様態は、停止構造の場において近接した領域を占める。両者はともに「環境との直接相互干渉によって停止を維持する」プロファイルである。

ただし、この構造的同位は、粘菌に認知やアニミズム的信念を帰属するものでは断じてない(NG6)。両者を隔てる決定的な差異は $\mu_\theta$ にある。粘菌は $\mu_\theta$ を内発的に持たず、環境条件に受動的に従って停止する。アニミズムを運用する主体は $\mu_\theta$ を内発的に持ち、応答チャネルの注入先を能動的に選び、組み替え、共有する。同じ「環境相互干渉型 $L_E$」というプロファイル領域にありながら、$\mu_\theta$ の内発性という軸において両者は明確に隔たっている。

これは「停止構造の場」が連続でありながら一様ではないことの好例である(\S{}0.4 NG6)。場は粘菌からアニミズムまで連続的に変化するが、$\mu_\theta$ という軸上で両者は異なる位置を占める。連続性は同一性ではない。

\bigskip


\subsection{次章への橋渡し}


アニミズムは、$L_E$ を環境に分散させて維持する形態である。停止を委ねる先が個々の自然物に分散しているため、$\Phi_{\text{net}}$ は分散的である。

この分散性には動力学的な帰結がある。分散した応答チャネルは、互いに競合・矛盾しうる。ある対象の応答と別の対象の応答が食い違うとき、停止構造に曖昧性が生じる。分散した $L_E$ ノードの間で停止判断が一致しない場合、推論主体はどの停止を採用すべきか決定できない。

次章 \S{}7 で扱う「神」は、この分散性の問題に対する動力学的応答として理解できる。分散した応答チャネルを単一の中心化された参照軌跡に統合することで、停止の曖昧性を低減する形態が神である。アニミズムから神への移行は、信念の進歩でも退歩でもなく、$\Phi_{\text{net}}$ の構造が分散から中心化へ変化する動力学的再編成である。

\bigskip


\subsection{要約}


本章で示したのは次の点である:

\begin{enumerate}
\item アニミズムは沈黙領域に応答チャネルを注入する操作である(定義 6.1)。$\rho_{\text{resp}}$ の局所的増大
\item アニミズムが行うのは事実主張ではなく停止構造の変換である
\item アニミズムは因果推論を否定せず保留する(説明遅延であって拒否ではない)
\item アニミズムは低コストの停止装置であり、その普遍性は構造的コストの低さによる
\item 応答的世界は責任の場を提供し、主体の行動可能性を支える
\item アニミズムは $L_E$ を環境直接相互干渉で維持する形態であり、$\Phi_{\text{net}}$ は環境的・分散的である
\item 粘菌の環境相互干渉型 $L_E$ と構造的同位にあるが、$\mu_\theta$ の内発性において明確に隔たる(連続だが一様でない)
\item アニミズムの分散性は停止の曖昧性を生み、これが \S{}7 の神(中心化)への動力学的動機となる
\end{enumerate}


次章 \S{}7 では、分散した $\Phi_{\text{net}}$ が中心化される形態——神——を分析する。

\bigskip
